\newpage
\section{Architettura di Sistema}

L'architettura software è stata ingegnerizzata seguendo un paradigma \textbf{stratificato} (\textit{layered architecture}), finalizzato a massimizzare la portabilità del codice e a garantire un netto disaccoppiamento tra l'astrazione dell'hardware e la logica applicativa. Tale approccio permette di gestire la complessità derivante dalla coesistenza di task asincroni e requisiti real-time tipici del SoC RP2040.


% --- SOTTOSEZIONE 1 ---
\subsection{Analisi dei Livelli Architetturali}

La separazione logica introdotta non è solo organizzativa, ma funzionale all'ottimizzazione del \textbf{determinismo} e della manutenibilità del sistema.

\subsubsection*{Livello Driver (Hardware Abstraction Layer)}
Rappresenta lo strato inferiore (\texttt{drivers/}) dove risiedono i componenti software che dialogano direttamente con i registri del SoC. Gestisce il condizionamento dei segnali e il campionamento tramite protocolli standard.
\begin{itemize}
    \item \textbf{I/O \& Bus:} Gestione I2C per sensori ambientali (\texttt{sensors\_i2c.h}) e display OLED (\texttt{display\_ssd1306.h}).
    \item \textbf{Timing:} L'uso del sottosistema \textbf{PIO} (\textit{Programmable I/O}) permette di sollevare la CPU dai compiti di timing critico per i segnali asincroni.
    \item \textbf{Comunicazione:} Gestione della telemetria radio tramite UART dedicata (\texttt{radio\_uart.h}).
\end{itemize}

\subsubsection*{Livello Core (Middleware \& System Management)}
Agisce come orchestratore centrale e ponte tra l'hardware e l'interfaccia utente (\texttt{core/}). Implementa i servizi critici di sistema:
\begin{itemize}
    \item \textbf{Storage:} Gestione resiliente dei dati tramite il filesystem \textit{LittleFS} e persistenza su Flash (\texttt{storage.h}).
    \item \textbf{Power \& Telemetry:} Monitoraggio dei consumi e pacchettizzazione dei dati di telemetria (\texttt{telemetry.h}).
    \item \textbf{System Manager:} Gestione degli errori (\texttt{error\_codes.h}) e coordinamento del boot del sistema.
\end{itemize}

\subsubsection*{Livello Application (User Interface \& Views)}
Costituisce l'interfaccia di alto livello (\texttt{ui/}). Gestisce il rendering grafico e la navigazione tra le schermate tramite un approccio modulare basato su \textbf{Views} (\texttt{view\_definitions.h}). Questo strato è agnostico rispetto all'origine fisica dei dati, interagendo esclusivamente con i gestori del livello Core tramite l'\texttt{UIManager}.

% --- SOTTOSEZIONE 2 ---
\subsection{Integrazione e Flusso Dati}

L'interazione tra i livelli avviene tramite interfacce asincrone. Questa granularità architettonica garantisce che le operazioni a latenza variabile, come i cicli di scrittura su Flash gestiti da \textit{LittleFS}, non interferiscano con la precisione temporale del parsing GNSS (affidato alla libreria esterna \textit{minmea}) o con la reattività del campionamento degli ingressi (\texttt{inputs.h}).
